\chapter{Kesimpulan dan Saran Pengembangan}
\label{ch:kesimpulan}

\section{Kesimpulan}
\begin{enumerate}
    \item Kami telah merancang \algname, sebuah \textit{block cipher} 256 bit berstruktur Feistel 18 ronde yang memadukan Conway's Game of Life dengan komponen SPN gaya AES. GoL dipakai di dua tempat: fungsi ronde $f_{\text{GoL}}$ pada ronde pertama dan terakhir, serta penjadwalan kunci yang menurunkan 18 \textit{round key} dari evolusi grid $16 \times 16$ dengan injeksi ulang \textit{nonce} dan \textit{whitening}. Struktur Feistel memungkinkan GoL, yang tidak dapat dibalik, dipakai sebagai fungsi ronde tanpa mengganggu proses dekripsi.
    \item Implementasi Python berhasil mengenkripsi dan mendekripsi berkas teks maupun biner pada kelima mode operasi (ECB, CBC, CFB, OFB, CTR). Seluruh hasil dekripsi identik bit demi bit dengan berkas asli, dibuktikan dengan kecocokan \textit{hash} SHA-256. \textit{Padding} PKCS\#7 menangani semua panjang pesan dengan benar, termasuk pesan kosong dan pesan yang panjangnya kelipatan ukuran blok.
    \item Secara statistik, \algname{} berperilaku seperti permutasi acak. \textit{Avalanche} terhadap \textit{plaintext} dan kunci masing-masing 49,95% dan 50,46%, difusi penuh tercapai setelah 4 dari 18 ronde, entropi \textit{ciphertext} sesuai dengan nilai harapan data acak pada setiap ukuran berkas, dan histogram \textit{ciphertext} rata, bahkan untuk \textit{plaintext} berupa byte nol. Kegagalan uji $\chi^2$ yang sesekali muncul terbukti sebagai fluktuasi acak melalui pengujian berulang.
    \item Analisis komponen menghasilkan tiga iterasi desain. Kotak-S afin pada versi awal meloloskan selisih \code{0x80} dengan peluang 1, sehingga diganti dengan Kotak-S AES. Pelipatan kunci dengan XOR membuat kunci $X \Vert X$ setara dengan kunci nol, sehingga diganti dengan SHA-256. \textit{Round key} yang diambil langsung dari grid hasil evolusi hanya memiliki proporsi bit 1 sebesar 0,362, sehingga ditambahkan \textit{whitening} yang menaikkannya menjadi 0,499. Temuan-temuan ini menunjukkan bahwa uji statistik seperti \textit{avalanche} dan entropi tidak cukup untuk menilai keamanan, karena versi awal pun lolos uji-uji tersebut.
    \item Beberapa keterbatasan masih tersisa pada versi akhir: \textit{ciphertext} tidak diautentikasi sehingga dapat dimanipulasi, batas peluang diferensial yang dapat dibuktikan baru $2^{-60}$, masih ada bias kecil yang bergantung pada posisi bit \textit{round key} (setara dengan kehilangan entropi sekitar 0,007 bit), dan kinerja implementasi Python hanya sekitar 60--70 KB/s.
\end{enumerate}

\section{Saran Pengembangan (Future Works)}
\label{sec:saran}
\begin{enumerate}
    \item \textbf{Autentikasi.} \textit{Encrypt-then-MAC} dengan HMAC-SHA256 atas $N \Vert C$ perlu ditambahkan, atau \algname{} dapat dikembangkan menjadi mode AEAD. Pesan galat \textit{padding} juga perlu disatukan agar tidak menjadi \textit{padding oracle}.
    \item \textbf{KDF untuk kata sandi.} SHA-256 sudah menghilangkan kunci ekuivalen, tetapi SHA-256 adalah fungsi \textit{hash} yang cepat. Untuk kunci yang berasal dari kata sandi, sebaiknya disediakan KDF lambat dengan \textit{salt}, seperti PBKDF2 atau Argon2, agar serangan tebakan kata sandi menjadi mahal.
    \item \textbf{Ekstraksi \textit{round key} yang lebih seragam.} Bias kecil per posisi bit dapat dihilangkan, misalnya dengan melewatkan grid hasil \textit{whitening} ke Kotak-S atau fungsi \textit{hash} sebelum dilipat, atau dengan mengacak urutan sel pada pengisian spiral.
    \item \textbf{Analisis formal.} Batas jumlah minimum Kotak-S aktif perlu dihitung dengan MILP untuk kriptanalisis diferensial dan linear. Kajian terhadap serangan \textit{related-key} juga perlu dilakukan, karena \textit{nonce} publik ikut menentukan \textit{round key}. Keluaran \textit{cipher} sebaiknya juga diuji dengan rangkaian uji statistik NIST SP 800-22 \cite{rukhin2010}.
    \item \textbf{Kinerja dan keamanan implementasi.} Laju sekitar 60--70 KB/s dapat ditingkatkan dengan implementasi bit-sliced dalam C/Rust atau dengan memproses blok secara paralel pada mode ECB dan CTR. Implementasi \textit{constant-time} juga diperlukan, karena \textit{lookup} Kotak-S dan rotasi dengan bilangan bulat Python berpotensi membocorkan informasi melalui waktu eksekusi.
    \item \textbf{Perbaikan antarmuka.} Mode operasi dapat disimpan di \textit{header ciphertext} agar penerima tidak perlu mengetahuinya terlebih dahulu. Galat validasi kunci sebaiknya dilempar sebagai \textit{exception}, bukan hanya dicetak. \textit{Counter} CTR juga perlu dibungkus modulo $2^{256}$ untuk menghindari \textit{overflow} pada \textit{nonce} yang sangat besar, walaupun peluang kejadiannya dapat diabaikan.
\end{enumerate}
